\documentclass[11pt,reqno]{amsart}
\usepackage{../notes/math327-lecture}
\newcommand{\Cl}{\operatorname{Cl}}
\newcommand{\F}{\mathbb{F}}
\newcommand{\C}{\mathbb{C}}
\newcommand{\SO}{\operatorname{SO}}
\pagestyle{plain}
\setlength{\parskip}{0.18em}
\setlist[itemize]{leftmargin=1.3em,topsep=0.1em,itemsep=0.12em,parsep=0pt}
\hypersetup{pdftitle={MATH 327 Midterm 2 Review Sheet},pdfauthor={Manish Patnaik},pdfsubject={Fall 2026 Midterm 2 coverage}}
\begin{document}
\begin{center}
\colorbox{lightblue}{\begin{minipage}{0.94\textwidth}
\vspace{0.5em}
{\color{courseblue}\bfseries\Large Midterm 2 Review Sheet\hfill MATH 327}\\[0.35em]
Algebra 1 \hfill Fall 2026\\
Instructor: Manish Patnaik
\vspace{0.5em}
\end{minipage}}
\end{center}
\textbf{Date:} Friday, November 6, 2026, in class.\quad
\textbf{Format:} Closed book, 50 minutes.\\
\textbf{Weight:} 25\%.\quad
\textbf{Coverage:} Lectures 15--23 (through Monday, November 2).

\smallskip
\noindent\textbf{Textbook references:} M. Artin, \emph{Algebra}, 2nd edition.
The sections below provide background for the listed topics; the lecture notes
and this sheet determine exam coverage.

\section*{Topics}
\subsection*{1. The correspondence theorem and \texorpdfstring{$\varphi:S_4\to S_3$}{the map from S4 to S3} (Lectures 15--16, 18)}
\noindent\textbf{Artin sections:} 2.10; 6.10.\par
\begin{itemize}
\item For a homomorphism $\varphi:G\to\mathcal G$ and $\mathcal H\leq\mathcal G$,
the preimage $\varphi^{-1}(\mathcal H)=\{g\in G:\varphi(g)\in\mathcal H\}$ is a
subgroup of $G$ containing $K=\ker\varphi$.
\item \textbf{Correspondence theorem:} when $\varphi:G\to\mathcal G$ is surjective
with kernel $K$, the maps $H\mapsto\varphi(H)$ and
$\mathcal H\mapsto\varphi^{-1}(\mathcal H)$ give a bijection between subgroups
$H\leq G$ containing $K$ and subgroups $\mathcal H\leq\mathcal G$.
Under this bijection, $H\trianglelefteq G$ if and only if
$\varphi(H)\trianglelefteq\mathcal G$, and $|H|=|K|\,|\varphi(H)|$ when $H$ is finite.
Review the proofs of bijectivity, normality in both directions, and the order formula.
\item \textbf{The map $\varphi:S_4\to S_3$:} $S_4$ acts on the three partitions of
$\{1,2,3,4\}$ into two unordered pairs,
$\Pi_1=\bigl\{\{1,2\},\{3,4\}\bigr\}$,
$\Pi_2=\bigl\{\{1,3\},\{2,4\}\bigr\}$, and
$\Pi_3=\bigl\{\{1,4\},\{2,3\}\bigr\}$.
Compute $\varphi(\sigma)$ for given $\sigma\in S_4$. Its kernel is the Klein four group
$K=V=\{1,(1\,2)(3\,4),(1\,3)(2\,4),(1\,4)(2\,3)\}$, a normal non-cyclic subgroup of order $4$.
\item \textbf{Subgroups of $S_4$ containing $K$:} the six subgroups of $S_3$ lift to
the six subgroups of $S_4$ containing $K$: $K$ (order $4$),
$A_4=\varphi^{-1}(\langle(1\,2\,3)\rangle)$ (order $12$, normal), three copies
$H_1,H_2,H_3$ of $D_8$ (order $8$, not normal), and $S_4$ (order $24$).
\end{itemize}

\subsection*{2. Group actions and dihedral groups (Lectures 17--18)}
\noindent\textbf{Artin sections:} 6.1, 6.4, 6.7, 6.11.\par
\begin{itemize}
\item \textbf{Left group actions:} a map $G\times S\to S$, $(g,s)\mapsto g\cdot s$,
satisfying $1\cdot s=s$ and $(gh)\cdot s=g\cdot(h\cdot s)$. Know the natural actions
of $S_n$ on $\{1,\ldots,n\}$ and $\GL_n(\R)$ on $\R^n$.
\item \textbf{Actions as homomorphisms:} an action of $G$ on a finite set
$S=\{s_1,\ldots,s_d\}$ induces $\rho:G\to S_d$ defined by $\rho(g)(i)=j$ when
$g\cdot s_i=s_j$. Its kernel is $\{g\in G:g\cdot s=s\text{ for all }s\in S\}$.
\item \textbf{Dihedral groups $D_{2n}$ ($n\geq3$):} generated by a rotation $x$ of
order $n$ and a reflection $y$ of order $2$ satisfying $yxy^{-1}=x^{-1}$
(equivalently $yx=x^{n-1}y$). Every element has a unique form $x^k$ or $x^ky$ with
$0\leq k<n$, so $|D_{2n}|=2n$.
\item \textbf{Symmetries of polygons:} work with $D_8$ (square) and $D_{12}$
(regular hexagon). Compose rotations and reflections right to left, identify the
reflection axes of the elements $x^ky$, and write the permutations induced on
vertices or diagonals in cycle notation.
\end{itemize}

\newpage
\subsection*{3. Direct products: external and internal (Lectures 18--19)}
\noindent\textbf{Artin sections:} 2.11.\par
\begin{itemize}
\item \textbf{External direct product $G_1\times G_2$:} componentwise multiplication,
identity $(e_1,e_2)$, inverses $(g_1^{-1},g_2^{-1})$, and
$|G_1\times G_2|=|G_1|\,|G_2|$ when finite. The factor subgroups
$\widetilde G_1=G_1\times\{e_2\}$ and $\widetilde G_2=\{e_1\}\times G_2$ are normal
in $G_1\times G_2$ and satisfy $\widetilde G_1\cap\widetilde G_2=\{(e_1,e_2)\}$.
\item \textbf{Products of cyclic groups:} if $\gcd(r,s)=1$, then
$C_{rs}\cong C_r\times C_s$. Review the proof using $\Phi(x^i)=(z^i,y^i)$ and why
coprimality is essential (for instance, $C_4\not\cong C_2\times C_2$).
\item \textbf{Product sets:} for $H,K\leq G$, the set $HK=\{hk:h\in H,\ k\in K\}$
need not be a subgroup (as in $S_3$). If $H\trianglelefteq G$ or
$K\trianglelefteq G$, then $HK=KH$ and $HK\leq G$.
\item \textbf{Internal direct products:} if $H\trianglelefteq G$, $K\trianglelefteq G$,
$H\cap K=\{1\}$, and $HK=G$, then $G\cong H\times K$ via $hk\mapsto(h,k)$.
Review unique factorization from $H\cap K=\{1\}$ and the commutator argument showing
$[h,k]=hkh^{-1}k^{-1}\in H\cap K=\{1\}$, so $H$ and $K$ commute (compare $S_3\cong C_3\rtimes C_2$).
\end{itemize}

\subsection*{4. Quotient groups and the first isomorphism theorem (Lectures 19--20)}
\noindent\textbf{Artin sections:} 2.12.\par
\begin{itemize}
\item \textbf{Quotient group $G/N$:} for $N\trianglelefteq G$, the coset product
$(g_1N)(g_2N)=(g_1g_2)N$ is well defined, with identity $N=1N$ and inverse
$(gN)^{-1}=g^{-1}N$. Know why normality is both sufficient and necessary for
well-definedness, and be able to show how the rule fails when $H\not\trianglelefteq G$.
\item \textbf{Quotient map:} $\pi:G\to G/N$, $g\mapsto gN$, is a surjective
homomorphism with $\ker\pi=N$; thus every normal subgroup is a kernel. For finite
$G$, $|G/N|=[G:N]=|G|/|N|$.
\item \textbf{First isomorphism theorem:} for a homomorphism $\varphi:G\to G'$ with
$N=\ker\varphi$, the map $\overline\varphi:G/N\to\im\varphi$ given by
$\overline\varphi(gN)=\varphi(g)$ is a well-defined isomorphism, so
$G/\ker\varphi\cong\im\varphi$. Review well-definedness, the homomorphism property,
injectivity, and surjectivity onto $\im\varphi$.
\item \textbf{Universal property of the quotient:} a homomorphism $f:G\to H$ factors
as $f=\overline f\circ\pi$ for a (unique) homomorphism $\overline f:G/N\to H$ if and
only if $N\subseteq\ker f$.
\item \textbf{Examples:} $\Z/n\Z$, $\GL_n(\R)/\SL_n(\R)\cong\R^\times$,
$S_n/A_n\cong\{\pm1\}$, $S_4/V\cong S_3$, $\C^\times/U\cong\R_{>0}$,
$(\R,+)/\Z\cong U$, and quotients of upper triangular matrix groups.
\end{itemize}

\newpage
\subsection*{5. Orbits, stabilizers, and the class equation (Lectures 21--22)}
\noindent\textbf{Artin sections:} 6.7--6.9; 7.2--7.3.\par
\begin{itemize}
\item \textbf{Orbits and stabilizers:} for $G\curvearrowright S$ and $s\in S$, the
orbit is $\mathcal O_s=G\cdot s$ and the stabilizer is
$G_s=\{g\in G:g\cdot s=s\}\leq G$. Orbits partition $S$; an action is transitive if
there is a single orbit. If $s'=a\cdot s$, then $G_{s'}=aG_sa^{-1}$.
\item \textbf{Orbit--stabilizer theorem:} the map $gG_s\mapsto g\cdot s$ is a
well-defined bijection $G/G_s\to\mathcal O_s$. When $G$ is finite,
$|G|=|\mathcal O_s|\,|G_s|$ and $|\mathcal O_s|=[G:G_s]$ divides $|G|$.
\item \textbf{Conjugation, centralizers, and center:} $G$ acts on itself by
$g\cdot x=gxg^{-1}$. The orbit of $x$ is its conjugacy class $\Cl_G(x)$ and the
stabilizer is its centralizer $C_G(x)=\{g\in G:gx=xg\}\leq G$. The center is
$Z(G)=\bigcap_{x\in G}C_G(x)$, and $x\in Z(G)\iff C_G(x)=G\iff\Cl_G(x)=\{x\}$.
\item \textbf{Class equation:} for finite $G$, every conjugacy class size divides $|G|$, and
$|G|=\sum_{i=1}^k|\Cl_G(x_i)|=|Z(G)|+\sum_{x_i\notin Z(G)}[G:C_G(x_i)]$.
\item \textbf{Applications:} class equation and centralizers of $S_3$ ($6=1+3+2$);
if $G/Z(G)$ is cyclic then $G$ is abelian; a group of prime-power order $|G|=p^n$
($n\geq1$) has $Z(G)\neq\{1\}$, and every group of order $p^2$ is abelian;
counting $|\GL_2(\F_3)|=48$ and $|\SL_2(\F_3)|=24$.
\end{itemize}

\subsection*{6. Conjugacy in \texorpdfstring{$S_n$ and $A_n$, and simplicity of $A_5$}{Sn and An, and simplicity of A5} (Lecture 23)}
\noindent\textbf{Artin sections:} 7.4--7.5.\par
\begin{itemize}
\item \textbf{Normal subgroups and conjugacy classes:} $N\leq G$ is normal if and
only if $N$ is a union of conjugacy classes of $G$. For finite $G$, $|N|$ must both
divide $|G|$ and be a sum of class sizes including $1$.
\item \textbf{Cycle types in $S_n$:} an $r$-cycle has order $r$ and sign $(-1)^{r-1}$;
a product of disjoint cycles of lengths $r_1,\ldots,r_k$ has order
$\operatorname{lcm}(r_1,\ldots,r_k)$. Two permutations in $S_n$ are conjugate if and
only if they have the same cycle type. Count class sizes and write the class
equations of $S_3$, $S_4$ ($24=1+6+8+6+3$), and $S_5$ ($120=1+10+20+30+24+15+20$).
\item \textbf{Splitting in $A_n$:} for $\sigma\in A_n$,
$C_{A_n}(\sigma)=C_{S_n}(\sigma)\cap A_n$. If $C_{S_n}(\sigma)$ contains an odd
permutation, then $\Cl_{A_n}(\sigma)=\Cl_{S_n}(\sigma)$. If
$C_{S_n}(\sigma)\subseteq A_n$, then $\Cl_{S_n}(\sigma)$ splits into two
$A_n$-classes of equal size $\tfrac12|\Cl_{S_n}(\sigma)|$.
\item \textbf{Simplicity of $A_5$:} in $A_5$, the $5$-cycles split into two classes of
size $12$ represented by $(1\,2\,3\,4\,5)$ and $(1\,2\,3\,5\,4)$, while the
$3$-cycles ($20$) and double transpositions ($15$) do not split, giving the class
equation $60=1+20+15+12+12$. Review the proof that no sum of these class sizes
(including $1$) is a proper divisor of $60$, so $A_5$ is simple.
\end{itemize}

%\subsection*{7. Rotations in $\R^3$ and poles (Lecture 24)}
%\begin{itemize}
%\item $\SO_3(\R)=\{A\in\GL_3(\R):A^{\mathsf T}A=I,\ \det A=1\}$ preserves dot products.
%Every $A\in\SO_3(\R)$ has $1$ as an eigenvalue; Euler's rotation theorem states that
%elements of $\SO_3(\R)$ are precisely the rotations of $\R^3$.
%\item For a finite subgroup $G\leq\SO_3(\R)$ of order $N$, $G$ acts on its set of poles
%$\mathcal P\subseteq S^2$, and point stabilizers $G_p$ are cyclic.
%\item \textbf{Pole-counting formula:} $\sum_{p\in\mathcal P}(|G_p|-1)=2(N-1)$, or
%$\sum_{i=1}^k(1-1/r_i)=2-2/N$, leading to two or three pole orbits and the classification
%patterns for finite rotation groups.
%\end{itemize}


\end{document}
